\documentclass[11pt,a4paper]{article}
% Compile with XeLaTeX (fontspec and xeCJK).
\usepackage{fontspec}
\setmainfont{LibertinusSerif}[Extension=.otf,UprightFont=*-Regular,BoldFont=*-Bold,ItalicFont=*-Italic,BoldItalicFont=*-BoldItalic]
\setsansfont{LibertinusSans}[Extension=.otf,UprightFont=*-Regular,BoldFont=*-Bold,ItalicFont=*-Italic]
\usepackage{xeCJK}
\setCJKmainfont{Hiragino Mincho ProN}
\xeCJKDeclareCharClass{Default}{"00B0, "00B7, "2013, "2014, "2018 -> "201F, "2026}
\usepackage[english]{babel}
\usepackage[a4paper,margin=2.5cm,top=2.3cm]{geometry}
\usepackage{graphicx}
\usepackage{xcolor}
\usepackage{titlesec}
\usepackage{fancyhdr}
\usepackage[hidelinks]{hyperref}
\definecolor{accent}{RGB}{31,73,125}
\definecolor{muted}{RGB}{93,104,119}
\hypersetup{colorlinks=true,urlcolor=accent,linkcolor=accent,pdfauthor={Romain Fernex},pdftitle={Japanese women in Meiji-era Japan : des trajectoires féministes à l’épreuve du conservatisme}}
\setcounter{secnumdepth}{0}
\titleformat{\section}{\large\bfseries\color{accent}}{}{0em}{}
\setlength{\parindent}{1.2em}
\setlength{\parskip}{0.35em}
\pagestyle{fancy}
\fancyhf{}
\fancyhead[L]{\small\color{muted}Romain Fernex}
\fancyhead[R]{\small\color{muted}\itshape Japanese women in Meiji-era Japan}
\fancyfoot[C]{\small\color{muted}\thepage}
\renewcommand{\headrulewidth}{0pt}
\newcommand{\partbreak}{\par\begin{center}\color{muted}$\ast$\quad$\ast$\quad$\ast$\end{center}\par}

\begin{document}
\thispagestyle{empty}
\noindent{\small\sffamily\color{accent}\MakeUppercase{History of Japan}}\par\medskip
\noindent{\LARGE\bfseries Japanese women in Meiji-era Japan\par}
\smallskip\noindent{\large\itshape Feminist trajectories put to the test by conservatism\par}
\medskip
\noindent{\color{accent}\rule{\linewidth}{0.8pt}}\par\smallskip
\noindent{\small\color{muted}Romain Fernex · 2023 · Article in English}\par\bigskip

On 8 July 1853, United States Commodore Matthew Perry entered Tokyo harbour aboard the frigate USS Susquehanna to demand Japan's reopening to international trade. Only fifteen years later, Japan's Sakoku policy (a policy of very restricted openness to the outside world) was officially lifted, and Emperor Komei's death marked the transition from the Tokugawa to the Meiji era. The latter was distinguished by major changes in every sphere, particularly industry and politics, with the end of the shogunate. Yet, in the world's eyes, this revolution seemed to have passed one part of Japan's population by: women. Whatever their status, they appeared neglected by this transformation, condemned to remain prisoners of an age dominated by patrilineal descent and Confucian values that drastically limited their freedoms both within the family and in public life. Does this mean that the Meiji era changed nothing for women in Japan? Advances in sociology and historical research now allow us to say quite the contrary. Long rendered invisible, Japanese women's struggle against a deeply patriarchal society did not wait until the 1960s, or even the Taisho era, to organise itself and develop its ideological structure and demands. Indeed, traces of some women's involvement in politics can already be observed during the Tokugawa era, although they tended to remain isolated. Why, then, choose the Meiji era? How did it mark a new beginning in Japanese women's struggle for independence? To answer this question, this article will analyse the birth and trajectories not of feminism but of the feminisms that emerged at the time, presenting their major participants. Here, we understand feminism as ``an awareness of the constraints placed upon women because of their sex and attempts to remove those constraints and establish a more equitable gender system involving new roles for women and new relations between women and men''{[}1{]}. We will therefore consider as feminist any individual, man or woman, aware of this inequality and speaking out against it. Finally, although using the term ``feminism'' here is not strictly an anachronism, it was very rarely, if ever, used by the activists discussed in this article. The concept was still recent, and invoking it entailed considerable risk in a conservative setting, in Japan as elsewhere.

To address this question, we will follow a chronological order, first discussing the transition from the Tokugawa to the Meiji era and the redefinition of women's role amid political unrest and intellectual dynamism, partly due to a new openness towards the West. We will then see how these new feminist movements responded to the reassertion of a conservative state intent on maintaining control over a changing society between 1884 and 1910. Finally, we will examine the end of the Meiji era and the culmination and transmission of a feminist body of thought that outlived these initial movements to enrich feminist struggles in subsequent eras.

\partbreak

First, we must assess Japanese women's situation during the Tokugawa era. Its inflexible, authoritarian political regime was distinguished by particularly strong state control over both private and public life. Japanese women, especially samurai women, who nevertheless occupied the upper ranks of society, were reduced to their reproductive function, with no rights over either husbands or children. Women of lower social origins faced somewhat fewer family constraints but remained exploited in fields or factories, paying with their time for the little freedom they enjoyed compared with their wealthier counterparts. The prevalence of Confucian values particularly hostile to their emancipation is visible in important contemporary texts such as the Onna Daigaku (Greater Learning for Women), which sought to describe the ideal Tokugawa woman. Women were thus attributed many faults, such as laziness or disobedience, supposedly making them slaves to their husbands, the only presumed respectable person in the household. Similarly, the principle of Danson Johi (respect men, despise women) dominated many aspects of society, making gender equality an exceedingly distant concept. But did this model go unchallenged? The American historian Anne Walthall identified women's presence in Tokugawa political life through, among other things, numerous signatures on petitions submitted by villages against imperial authority. This demonstrates women's involvement alongside men in legal struggles against repressive authorities. Moreover, several women, such as Matsuo Taseko, a village chief's wife, managed at least partly to transcend the social conventions bearing down on them and establish a real public presence. Yet, as Walthall notes, these manifestations remained sporadic and never produced a convergence of struggles aimed at improving Japanese women's status in its broadest sense.

The shock of some Japanese men and women's encounters with Western women was therefore all the greater. At the end of the Tokugawa era, during a Japanese government diplomatic mission to the White House, the presence and behaviour of US President James Buchanan's niece surprised the Japanese envoys. She presided over the occasion and proved more than capable of asserting herself at the table, something entirely unthinkable for a Japanese woman given the status society imposed on her. This encounter, one among many, helped a Japanese government undergoing reform realise that women's status was one of the elements radically separating the Western world from the Japanese archipelago. Japan was still traumatised by the unequal treaties imposed by the United States and sought more urgently than ever to overcome its developmental lag behind the West to regain independence. A new discourse therefore emerged among certain Japanese elites: since women's education and status seemed important to Western success, Japan must change in this respect too. This revelation produced a particularly ambiguous policy oscillating between ostentatious measures to please the West, progressive rhetoric and resurgent conservatism in response to what was perceived as a challenge to Japan's traditional social model.

Education was one of the first pillars of this reform. As early as 1872, the American David Murray, advising the Japanese government, called for women to be included in education reform and for women's schools to be established. This call did not come solely from abroad, as the engagement of certain Japanese elites in the press demonstrates. The Meirokusha (Meiji Six Society), founded in 1873 and particularly influential at the time, published several articles on women's rights in its journal, the Meiroku Zasshi. Several contributors stood out for their critical positions on women's treatment in Japanese society, considered a genuine progressive commitment at the time. Thus, Viscount Mori Arinori, a founding member and Japan's first ambassador to the United States, championed women's right to education through the principle of Ryosai Kenbo (good wife, wise mother). Although he clearly avoided calling for equality between men and women, the assertion that women had a real role in the family as educators, and therefore needed training to fulfil it, was particularly innovative. Following an initial wave of schools for samurai girls in the 1850s, the government introduced an ambitious education policy in 1872 with the publication of the Gakusei (Fundamental Code of Education), advocating harmonised subjects regardless of gender or social class and declaring education a right for everyone. This accompanied the establishment of many private girls' schools, including Ueda Jogakko in Tokyo in 1872 and Tokyo Joshi Shihan Gakko (Tokyo Women's Normal School) in 1874, with the empress's support. Nevertheless, it had limitations: although schooling was now a universal right, few families were prepared to fund their daughters' studies, particularly since much of society still viewed girls' education negatively. Moreover, girls enrolling in these schools were often forced to leave before graduation for family or financial reasons. However, the government was not the only participant in this renewal of girls' education. Westerners also played an important role by establishing missionary schools teaching young Japanese men and women the basics of Christianity{[}1{]}, but also physics, mathematics and Western literature. Ferris Jogakuin (Ferris Girls' School), founded by the missionary Clara Hepburn in 1859 and later taken over by Mary Kidder, was one of many. These formed an extensive network of institutions open to Japanese women regardless of social class. This network, alongside that developing in certain government schools, was particularly important to the emergence of a genuine feminist consciousness. Young women came together, often far from their families, and encountered Western works frequently more open to questions of women's place{[}2{]}. This new environment, shaped by education's often formative years, gave rise to new forms of socialisation for Japanese women and, for some, entry into politics, bringing issues conceived by women and directly affecting them into the public sphere.

This political and social novelty fitted into a framework of social demands and the politicisation of struggles against the government. In 1877, the last pocket of military resistance to the Meiji government was definitively extinguished, leaving reformers only the political route. Indeed, from 1873, several politicians, including former samurai Itagaki Taisuke, withdrew from government to found a new reform movement, the Freedom and People's Rights Movement (Jiyu Minken Undo), seeking genuine democracy. Initially hostile to admitting women, it gradually recognised their importance as a political force in their own right, and several members stood out for particularly open positions on women's rights. Several Western-inspired leaders, such as Ueki Emori, saw improving women's status, albeit primarily within the home, as a genuine driver of Japan's development. The movement therefore gradually offered platforms to women, particularly on questions of women's status. In 1881, the young Japanese woman Kishida Toshiko gave her first public speech to considerable acclaim. Exceptionally educated, having tutored the empress for several years, she opposed the system of concubinage{[}3{]} prevailing both at court and in Japanese households, and demanded women's assertion in both private and public life. Coming from Kyoto, she was particularly familiar with young Japanese women's precarious status{[}4{]}, and her knowledge of leading Western feminists, such as Millicent Fawcett and the suffragette movement, gave her words particular resonance. In October 1883, she organised the first meeting of the Joshi Daienzetsukai (Kyoto Women's Reading Society), attracting nearly 2,000 women. She also published extensively, criticising the government and defending Japanese women's status by invoking great female figures from Japanese history, such as Murasaki Shikibu, author of Genji Monogatari{[}5{]} (The Tale of Genji), or Western women such as Elizabeth Browning. Nevertheless, this proliferation of women's discussion groups and magazines devoted to women's rights soon confronted an authoritarian government worried by an increasingly powerful hostile discourse perceived as threatening Japan's traditional model. Growing repression and divisions led to the Liberal Party's breakup in 1884. The first party of the Meiji era, born of the Jiyu Minken Undo, its collapse dealt a severe blow to the young feminist movements, most of which had been closely connected to it and found in it a real means of political expression.

1884 was indeed a major year in many respects for Japanese society and particularly for women's liberation movements. Despite tightening censorship and worsening economic and political conditions, it was a year of maturation and expansion for Japanese feminist movements. Mary Leavitt, a missionary and representative of the WCTU (Woman's Christian Temperance Union), arrived in Japan that year and delivered her first public speech in Tokyo just two years later. At this event, sponsored by the Jogaku Zasshi{[}6{]}, one of the leading contemporary journals concerned with women's condition, she addressed an exclusively female audience, defending women's role while denouncing vices obstructing their development. She firmly condemned Japanese men's use of alcohol and tobacco and criticised concubinage, still prevalent despite several cosmetic measures. Although much of the audience recognised her remarks' religious subtext, many heard in them the problems they encountered daily. Mary Leavitt also promoted women's education through concrete examples of European women just beginning to enter professions previously reserved for men, such as medicine and law. This speech came as certain Japanese feminists, including Toshiko Kishida, became aware of women's actually difficult circumstances across the Atlantic. Far from seeing grounds for despair, Kishida argued in an 1884 article that feminist struggle offered a means of finally closing the gap with the West by achieving what it had not yet accomplished. At this time, young women inspired by missionaries' speeches and women such as Kishida also entered the struggle for women's rights, motivated by Christianity, socialism, or both.

Fukuda Hideko was one of them. Still a teenager when Kishida spoke in 1882, she drew from it a commitment to the cause that she would retain throughout her life. After several initiatives, including a private girls' school closed in 1884 for lack of funds, she was arrested by police during disturbances in Osaka in November 1885{[}7{]}. While she was imprisoned, in 1886, Yajima Kajiko, another feminist inspired by Mary Leavitt's words that year, founded the Tokyo Fujin Kyofukai (Tokyo Women's Reform Society), which soon became important in spreading feminist ideas in Japan. Shortly after Fukuda's release, in 1889, the Meiji government adopted a first constitution particularly unfavourable to women's status. Through Article 5 of the Police Security Regulations, it denied women any right to participate in politics locally or nationally. Among other things, it prohibited women from attending future parliamentary meetings deciding new laws. This provoked an outcry in Japanese society, particularly among activists of the former Jiyu Minken Undo, to which some members of the new government belonged. Yajima Kajiko and the Tokyo Fujin Kyofukai launched numerous protests against the ban on women entering the Diet. This was one of the first nationwide campaigns organised by a women's rights movement, and its petition secured a very slight relaxation of the initial measure. While most Tokyo Fujin Kyofukai activities ordinarily condemned concubinage and prostitution networks{[}7{]}, Japanese women here considered their role as political participants in public life rather than solely within the private sphere of the home. In this particularly turbulent context, Fukuda Hideko founded another private school for young women and drew closer to the young socialist movement led by Sakai Toshihiko. This heterogeneous movement formed in 1901 in response to rising government authoritarianism, marked by the disastrous 1898 Civil Code reform reaffirming the restrictive view of women that had prevailed during the Tokugawa era. Coming after the 1895 war between China and Japan, this reform showed that Meiji authoritarianism continued to grow under pressure, with the return to traditional values expressing a desire for renewed stability. The socialist movement therefore united Christians, former Jiyu Minken Undo members, activists and others who shared a belief in socialism as a path to Japan's desired development and, above all, were deeply critical of the government's direction.

This movement also echoed another major transformation beginning shortly before the 1880s: the massive industrialisation of Japan. To modernise its economy, Japan turned towards silk production and export, opening numerous mills and factories. From 1894 to 1912, women accounted for 60\% of industrial workers in Japan. Many young women from poor families were forcibly recruited into workplaces where conditions were particularly harsh, as portrayed in Satsuo Yamamoto's 1979 film Nomugi Pass. Workers received only meagre pay, immediately passed to their families, preventing even a hint of financial independence despite their efforts. This difficult reality nevertheless generated important local protests, particularly major strikes at certain workplaces. In June 1886 in Amamiya, female workers, soon joined by their few male colleagues in the mills, protested wage cuts and demanded higher bonuses, often their only real source of income. The mobilisation, the first of this scale in the Meiji era, lasted several weeks and ended in victory for the strikers. Although the first of a long series, not all succeeded, and no national movement emerged because the dormitory system isolated female workers. Moreover, differences in social background between these workers and most feminist activists, largely middle-class and possessing at least some education, severely limited communication between the two groups.

What, then, of the socialist feminist movement, combining rupture and continuity with the many often Christian-inflected upper-middle-class women's movements, such as Toshiko Kishida's? Several members had already participated in social struggles by promoting trade-union formation during the 1897 economic crisis. Women joining the movement felt the need to prioritise more issues concerning women's condition, prompting them to form a distinct group focused on these problems. Some leaders supported this effort by requiring every man to be accompanied by a woman at meetings to maintain a form of parity, and by allowing several women to speak publicly. However, the association dissolved in 1904 amid mounting government pressure during the war with Russia and internal divisions between Christians, materialists and other currents. Socialist feminism truly ended, however, after the ``Red Flag Incident'' of 22 June 1908. The government's release of anarchist activist Koken Yamaguchi prompted a substantial socialist mobilisation, welcoming the news with an unauthorised demonstration in Tokyo. Several participants, including Kanno Suga, ardently supported Yamaguchi's ideas and carried red flags bearing anarchist slogans. The mobilisation provoked intense police repression. The movement's magazine was banned, and barely a year later Kanno Suga was arrested for planning the emperor's assassination. She was executed in 1910 alongside twenty-four other members, definitively marking the end of the socialist endeavour.

\partbreak

This wave of repression ended many feminist movements by making their activities extremely risky. Yet many activists did not abandon the struggle. Despite tighter controls over missionary schools and the press, publications remained relatively regular, although often local and subject to censorship. In 1911, despite repression of opposition movements, the Seito association (better known as the Bluestockings) was founded. It was unique in having no initial political motivation. Its founder, Hiratsuka Raicho, the daughter of a wealthy liberal Tokyo family and educated at Naruse Jinzo's women's university, primarily wanted a platform to showcase women artists' work and initially had little awareness of problems concerning women's status. Becoming a successful female writer or artist remained particularly difficult apart from rare exceptions, and platforms available to women in art and politics were severely limited. From its first edition, the magazine received substantial support from established artists and financial assistance from the founder's mother. Naturally, celebrating female talent in any field provoked the anger of commentators who disparaged the initiative. Nevertheless, the magazine enjoyed considerable success among talented young women and women from diverse social classes. Among new members was Ito Noe, a young Japanese woman of humble origins who came to study in Tokyo after fleeing her native Kyushu following an arranged marriage. She advocated a more political direction, even at the risk of confrontation with the government. In April 1912, an issue critically examining Japan's family system resulted in a temporary government ban and more or less unofficial surveillance. Yet this did not prevent an increasingly political turn, with a complete series on women's problems a year later. The series proved highly successful with readers, and its content was relayed in major national newspapers. Topics included the international feminist movement through the suffragettes, the struggle for autonomy and equal rights, and women's sexuality. Fukuda Hideko also contacted the association, advocating socialism as a remedy, which elicited mixed reactions from some members.

Tensions increased further when the Home Ministry issued a list of women's magazines prohibited from publication later that year. Although the Bluestockings were not included, the association strengthened its commitment by changing its statutes. However, this increasingly political direction worried members who had joined primarily to gain visibility for their work. Dissensions grew until 1915, when an exhausted Raicho handed the association's leadership to Ito Noe. Isolated and with very limited funds, Ito Noe nevertheless kept it functioning for another year. That year proved particularly prolific, introducing important questions such as abortion and parenthood into public debate, questions that persisted long after the journal ended. Ito Noe's fate was tragic: arrested during the Tokyo earthquake in September 1923, she was killed in her cell alongside her nephew and husband by a radicalised police officer. Nevertheless, the Bluestockings left a powerful legacy, drawing, sometimes unconsciously, on preceding decades of struggle and laying foundations for the second wave of feminism during the Taisho era (officially beginning in 1912). Thus, in January 1920, Raicho, together with Ichikawa Fusae and Oku Mumeo, founded the Shinfujinkyokai (New Women's Society) and its associated magazine, Josei Domei (Women's League). Enriched by women's experience during the Meiji and early Taisho eras, this society became a major participant in feminist struggles over subsequent decades.

\partbreak

In conclusion, far from a period of stagnation and resignation for women, the Meiji era witnessed new forms of female socialisation and connections between previously disorganised feminist struggles. Although the political climate was clearly, overall, unfavourable to even modest reforms towards gender equality, this experience proved valuable for future movements, and many present-day Japanese feminist movements find solid ideological foundations there. Western contributions, combined with an emerging ``class consciousness'' and assisted by expanding media through newspapers' democratisation, thus produced a generation of feminist movements. They were admittedly diverse and sometimes critical of one another, yet united by a shared experience of being a woman in Japan and determined to restore women to their rightful place.

\textbf{Notes:}

\textbf{{[}1{]}} I use the definition proposed by the scholar Margot Badran in \emph{Feminists, Islam and Nation: Gender and the Making of Modern Egypt}. It is also used in Mara Patessio's work cited in the bibliography.

\textbf{{[}2{]}} This certainly does not mean that women's situation was radically better in the West or that contemporary Western texts were models of feminism. Yet they often offered a radically different perspective to young Japanese women previously exposed mainly to Confucian-influenced texts, primarily prescribing a morality that excluded them from any ambition to reflect on society. Works such as the English philosopher Herbert Spencer's \emph{Social Statics} and John Stuart Mill's \emph{The Subjection of Women}, published in 1851 and 1869 respectively, were therefore genuine revelations for some Japanese women who embraced the feminist cause.

\textbf{{[}3{]}} Polygamy was a recurring feature of Japanese households during the Tokugawa and Meiji eras. It involved concubines often treated similarly to the household head's wife and sometimes expected to provide him with an heir. Such an heir was generally considered legitimate, and the married woman often had to care for the child at her husband's request. Within Japanese society, this practice weakened wives' status and reinforced asymmetry within couples. Opposition to concubinage was therefore a recurring theme of feminist struggles concerning private life. Yet its form varied between groups and periods. Although attacks on concubinage initially targeted participating men and women indiscriminately, several feminists, such as Fukuda Hideko, developed more nuanced critiques recognising many concubines' particularly precarious status.

\textbf{{[}4{]}} Toshiko's Kyoto origins matter all the more because the Kyoto-Osaka region was the setting for particularly harsh practices by men. These included ``Daughters in boxes'', whereby marriageable girls were confined at home by their families until presentation to their future husbands. Kishida sharply criticised this radical denial of women's freedom in several speeches, including one of her most famous, on 12 October 1883 in Otsu. That speech led to her arrest and imprisonment for several days.

\textbf{{[}5{]}} Genji Monogatari is among the best-known works of Japanese literature both domestically and internationally. It was written in the 11th century, during the Heian era, by Murasaki Shikibu, an imperial court dignitary and one of the first poets to write in Japanese. Chinese had long dominated the imperial court, relegating Japanese to second place. Thus, many of the first successful Japanese-language authors were court women, demonstrating their learning and an aesthetic sensibility fully comparable to that of their male counterparts.

\textbf{{[}6{]}} Jogaku Zasshi was one of the Meiji era's widely circulated progressive journals. Publication began in July 1885 at the initiative of a man, Iwamoto Yoshiharu, an advocate of women's education and instructor at Meiji's Girls' School. It published numerous articles supporting the feminist cause, most written by women engaged in the struggle. Its influence lasted until its closure in 1904.

\textbf{{[}7{]}} Alongside concubinage, prostitution networks were a recurring theme in Japanese feminist struggles. The subject is particularly broad and undoubtedly deserves an article of its own, but some details are useful here. Society in the Tokugawa, Meiji and Taisho eras was marked by an institutionalised prostitution network sending numerous young Japanese women of limited means around the globe. They were forced to provide services to wealthy foreign clients or Japanese soldiers exhausted by imperial conquests. This ``industry'' became structured during the Meiji era through the Karayuki and Ameyuki-san phenomenon, developing in the late 19th century and persisting throughout the first half of the 20th. These structures laid the foundations for the deeply controversial ``comfort women'' system established by the Japanese empire in its conquered Asian territories during the Second World War.

\textbf{Bibliography:}

\begin{itemize}
\item
  Sharon L. Sievers, \emph{Flowers in Salt, The beginnings of Feminist consciousness in Modern Japan}, Stanford University Press, 1983
\item
  Mara Patessio, \emph{Women and Public life in Early Meiji Japan, The development of the feminist movement}, vol.71, University of Michigan Center for Japanese Studies, 2011
\item
  Anne Walthall, “Devoted Wives/Unruly Women: Invisible Presence in the History of Japanese Social Protest,” Signs 20.1 (1994): 106–36; “Edo Riots,” dans \emph{Edo and Paris: Urban Life and the State in the Early Modern Era}, édité par James L. McClain, John M. Merriman, et Ugawa Kaoru, 407–28 (Ithaca: Cornell University Press, 1994).
\item
  Margot Badran, \emph{Feminists, Islam and Nation~: Gender and the Making of Modern Egypt}, PUP, 1995, 19-20
\end{itemize}

\end{document}
